% Appendix A, protocol history and analysis plan (restructure plan, part 3). Sources: experiments/protocol.md (§1, §3,
% §5-§7, §15-§20); docs/decision_log.md (D15, D22, D32, D39, D40, D44, D48, D49, D51, D53-D61, D64, D71); the attempt
% records data/fresh/fresh-v1/attempts_CS-0*.json and the labelers' specs (frozen_at). Every number is a macro: \dv
% (evidence_ledger_design.md: each value checked against its source, each protocol hash recomputed), \fv, \plv.
% Moved here: protocol.tex's versions, the three amendments and the fresh stratum's timeline, its runs, its secondary
% measures and its reporting rules.
\section{Protocol history and analysis plan}
\label{app:protocol}\label{sec:protocol}\label{sec:amendments}
The protocol (\texttt{experiments/protocol.md}) was hashed and timestamped locally before the first evaluation run
(v0.1, \dv{protocol.v0.1.date}; D15); it is not a formal preregistration. Amendments v0.2 to v0.6
(\dv{protocol.v0.2.date} to \dv{protocol.v0.6.date}) served the construction pilot (D22, D32, D39, D40, D44). Later
ones are appended as new sections while the version line deliberately still reads v0.6, so each version from v0.6 on
is a byte prefix of the current file and its logged hash can be recomputed. Table~\ref{tab:amendments} lists every
version: what it fixes, what did not exist yet when it was hashed, and the relevant data that existed then. The
question of protocol~\S1 asks for the frequency with which omitted user corrections make a verifier reward an
unacceptable solution, and for the reconstruction and validation methods that reduce that error without rejecting valid
alternatives. This paper answers a conditional form of it (Section~\ref{sec:intro}): given a correction that adds or
retracts requirements, whether verifiers built without it accept violations of it, on controlled tasks. The design
cannot say how often corrections are lost in practice.

\IfFileExists{generated/design/amendments_table.tex}{\begingroup\scriptsize
\begin{longtable}{@{}p{0.12\linewidth}p{0.36\linewidth}p{0.21\linewidth}p{0.21\linewidth}@{}}
\caption{The protocol's versions (\texttt{experiments/protocol.md}), each hashed and timestamped locally (UTC; the first 8 hex digits of its sha256, recomputed from the git object of its commit or from a byte prefix of the current file, as the decision-log entry named records it). Fixed before: what did not exist yet when the version was hashed; after: the relevant data that existed then.}\label{tab:amendments}\\
\toprule Version & What it fixes & Fixed before & After \\\midrule
\endfirsthead
\toprule Version & What it fixes & Fixed before & After \\\midrule
\endhead
\midrule\multicolumn{4}{r}{continued} \\
\endfoot
\bottomrule
\endlastfoot
v0.1\newline \dv{protocol.v0.1.date}\newline \dv{protocol.v0.1.time}\newline \texttt{\dv{protocol.v0.1.hash}}, \dv{protocol.v0.1.entry} & The pilot's question, populations, unit of analysis, conditions, candidates and labels, metrics, comparisons and freezing (\S1--\S9). & Any evaluation run. & The controlled sessions; B1's construction run on CS-01's reference, used with the no-op and the pre-correction state as smoke checks. \\
v0.2 (\S10)\newline \dv{protocol.v0.2.date}\newline \dv{protocol.v0.2.time}\newline \texttt{\dv{protocol.v0.2.hash}}, \dv{protocol.v0.2.entry} & Youden's J and the families with an accepted negative; negatives by origin; the label rule D19; generator diversity as a limitation. & Any evaluation run. & The verifier-fidelity literature review. \\
v0.3 (\S11)\newline \dv{protocol.v0.3.date}\newline \dv{protocol.v0.3.time}\newline \texttt{\dv{protocol.v0.3.hash}}, \dv{protocol.v0.3.entry} & The final pilot run: harness 0.3.0 or later with every condition of every task re-run, B1 for CS-04 re-selected (D31), sanity grading, run conditions; harness reliability, contract recovery and independent construction checks as secondary analyses. & Any evaluation run under harness 0.3.0. & The superseded runs \texttt{pilot-v0.2} and \texttt{pilot-v0.2-h2} and the harness defects they showed (D28--D31). \\
v0.4 (\S12)\newline \dv{protocol.v0.4.date}\newline \dv{protocol.v0.4.time}\newline \texttt{\dv{protocol.v0.4.hash}}, \dv{protocol.v0.4.entry} & Harness 0.3.1 for a security defect; disclosures of duplicates, of B1's tests in the near-miss workspaces and of reference checks written with the pool; D19's labels re-proposed and every pair re-run. & The 0.3.1 label re-proposal and run (\texttt{pilot-v0.3.1}). & The internal review of \texttt{pilot-v0.3} (D36--D38). \\
v0.5 (\S13)\newline \dv{protocol.v0.5.date}\newline \dv{protocol.v0.5.time}\newline \texttt{\dv{protocol.v0.5.hash}}, \dv{protocol.v0.5.entry} & Harness 0.3.2 for a remaining defect; a further duplicate; corrected wording (what pre-registered means, the exceptions to generating the pool after freezing, the candidate recipe and the mapper instructions, reference checks written with the pool, generation costs not recorded); D19's labels re-proposed and every pair re-run. & The 0.3.2 label re-proposal and run (\texttt{pilot-v0.3.2}). & The dispositions review and the final claims audit of \texttt{pilot-v0.3.1}; the audit's completeness critic finished after the hash (D40). \\
v0.6 (\S14)\newline \dv{protocol.v0.6.date}\newline \dv{protocol.v0.6.time}\newline \texttt{\dv{protocol.v0.6.hash}}, \dv{protocol.v0.6.entry} & Harness 0.3.3 for a defect in the error path; an audit gap of the held-out generators and the reference-check writers; corrected wording; D19's labels re-proposed and every pair re-run (the fixes of D40--D43). & The 0.3.3 label re-proposal and the run \texttt{pilot-v0.3.3}. & A verification of the fixes by model agents (D44). \\
v0.7 (\S15)\newline \dv{protocol.v0.7.date}\newline \dv{protocol.v0.7.time}\newline \texttt{\dv{protocol.v0.7.hash}}, \dv{protocol.v0.7.entry} & Pilot families become development families, later pilot re-analyses exploratory; strata \texttt{fresh-v1} (variants, configurations, $k=\dv{design.k}$) and \texttt{natural-pr-v1} (S9 and S10 rules, family cap, stopping, size rule); the independent labeler; \dv{design.repeats.w} repeats per pair and a re-grade of the pilot; reporting rules. & Any generation for the new strata. & The pilot's results and the deterministic AIDev funnel with its S9 sample draw. \\
v0.8 (\S16)\newline \dv{protocol.v0.8.date}\newline \dv{protocol.v0.8.time}\newline \texttt{\dv{protocol.v0.8.hash}}, \dv{protocol.v0.8.entry} & Harness 0.3.5 for 0.3.4; path substitution and the audit for fresh attempts; identifiers; S6 corrected, S11 defined, S9 failures ledgered; reconstruction check for new tasks; reporting rules in the analysis code. & Any run, fresh attempt or natural verifier. & The first sample's S9 judgments (D50). \\
v0.9 (\S17)\newline \dv{protocol.v0.9.date}\newline \dv{protocol.v0.9.time}\newline \texttt{\dv{protocol.v0.9.hash}}, \dv{protocol.v0.9.entry} & S10 environments (offline sdist builds, \dv{screen.s10.wheelhouse}\,GB cap, closure rule), base-test subset, B1 unavailable, replacement of rejected families, storage. & Any S10 build, natural verifier or grading. & A read-only feasibility scout of the \dv{natural.scout.families} provisional families. \\
v0.10 (\S18)\newline \dv{protocol.v0.10.date}\newline \dv{protocol.v0.10.time}\newline \texttt{\dv{protocol.v0.10.hash}}, \dv{protocol.v0.10.entry} & Review rule for audit flags of fresh attempts; sensitivity view without every flagged attempt; harness-context disclosures. & Any grading or label. & CS-01's flags had been read (\dv{fresh.flags.read}; D58), while the last \dv{fresh.attempts.after.v0.10.ended} of the \dv{design.attempts} attempts ran (\dv{fresh.attempts.after.v0.10.started} started after the hash; the last ended at \dv{fresh.attempts.last}). \\
v0.11 (\S19)\newline \dv{protocol.v0.11.date}\newline \dv{protocol.v0.11.time}\newline \texttt{\dv{protocol.v0.11.hash}}, \dv{protocol.v0.11.entry} & S9 extended to the other \nat{s9.ext.comments} comments; families and size rule over both samples; order of work; bare-name labeler revisions. & Any extension judgment and any natural contract, verifier, held-out candidate or label. & S10 of the first sample (\nat{s10.first.passed} of \nat{s10.first.processed} families pass; D59). Post hoc. \\
v0.12 (\S20)\newline \dv{protocol.v0.12.date}\newline \dv{protocol.v0.12.time}\newline \texttt{\dv{protocol.v0.12.hash}}, \dv{protocol.v0.12.entry} & Each natural sample gets a run of its own: \texttt{natural-pr-v1} reports the first sample's families as case studies, \texttt{natural-pr-v2} the extension's (\S20.2, \S20.3). We made it to report the first sample's families before the extension is built, which needs agent usage beyond the budget (\S20.1). & Any P5 step and any natural held-out candidate. & P1--P4 of the first sample's families (D70). Post hoc, disclosed as v0.11 is. \\
\end{longtable}
\endgroup
}{\unfilled{amendments table (make\_tables.py -{}-design)}}

\paragraph{The amendments that bear on what the results can show.}
\begin{itemize}
\item v0.7 made the pilot's families development families, which the fresh-agent stratum, designed after the pilot's
  results were known, reuses with their frozen verifiers. It also fixed, before any grading of the new strata, the
  per-family table as the primary display and the oracle-gated view as secondary (protocol~\S15.7).
\item v0.10 fixed how the audit flags of fresh attempts are reviewed after the first family's flags had been read, but
  before any grading or label. Its preamble says it followed the \dv{design.attempts} attempts, which the attempt
  records contradict for \dv{fresh.attempts.after.v0.10.ended}; the sensitivity view does not depend on the review.
\item v0.11 extended the natural screening after the first replay yield was known (D59, D60): a post hoc design change,
  made before any natural contract, verifier, held-out candidate or label existed.
\item v0.12 gave each natural sample a run of its own (\S20). We made it to report the first sample's families before
  the extension is built, which needs agent usage beyond the budget (\S20.1). It was hashed after phases P1 to P4 of the
  first sample's families and before any P5 step or natural held-out candidate (D70, D71), and it is post hoc,
  disclosed as v0.11 is.
\end{itemize}

\paragraph{The fresh-agent stratum's order of work.} On \dv{fresh.date} (UTC):
\begin{itemize}
\item the independent labelers were frozen from \dv{fresh.labelers.first} to \dv{fresh.labelers.last} (D55);
\item the instructions and their hashes were logged at \dv{fresh.instructions.logged} (D56);
\item the attempts ran from \dv{fresh.attempts.first} to \dv{fresh.attempts.last} (D58); CS-01's audit flags were read at
  \dv{fresh.flags.read}, and protocol v0.10 was hashed at \dv{protocol.v0.10.time}, while the last
  \dv{fresh.attempts.after.v0.10.ended} attempts ran (\dv{fresh.attempts.after.v0.10.started} of them started after the
  hash);
\item D19's labels were proposed from \dv{fresh.labels.first} to \dv{fresh.labels.last}, and the transcripts of the
  labelers' writers were audited at \dv{fresh.audit.writers} (D58);
\item the run graded from \fv{run.started} to \fv{run.finished} (D61).
\end{itemize}
The natural arm's generation recipe was fixed, and its \dv{natural.templates.w} prompt templates hashed on
\dv{natural.templates.date} at \dv{natural.templates.time}, before any natural agent started (D64). Its builder
decisions await the owner's review (D65; Appendix~\ref{app:deviations}).

\paragraph{The natural run's order of work.} On \nrv{run.date} (UTC):
\begin{itemize}
\item D19's labels of the first sample's held-out candidates were proposed from \dv{natural.d19.first} to
  \dv{natural.d19.last}, as the run's first step, and the pre-flight passed (D73, D74, D76);
\item the run started at \nrv{run.started}; its first invocation was stopped before any pair was graded, and the run
  started again at \dv{natural.run.restarted} (D76);
\item the template of Section~\ref{sec:natural-cases} and the analysis fields were committed before any analysis of the
  run, while grading was still in progress (D77 records its state at \dv{natural.template.graded}); by then the build
  session's progress checks had shown it counts of outcomes over all graded pairs together, naming no verifier or
  candidate (D77);
\item the run ended at \nrv{run.finished}, and was then analysed in case-study mode and published, and its tables were
  generated (D78); the closing paragraph's first sentence and the restructure's other changes to the template followed
  (Section~\ref{sec:natural-cases}).
\end{itemize}

\paragraph{Strata and runs.} Each stratum has its own runs and analyses, never pooled without saying so
(protocol~\S15.2).
\begin{itemize}
\item \texttt{\PILOTRUN}, the construction pilot, graded every condition on the constructed pool under harness
  \plv{run.harness}, a container at a time, with no other sandbox workload on the host. Its analysis,
  planned in protocol v0.1 to v0.6, stands; since v0.7 it is a development study, whose later re-analyses are
  exploratory and reported apart (D49; Appendix~\ref{app:pilot}).
\item \texttt{\REGRADERUN} re-graded the same pairs under harness \fv{rg.harness} with \fv{rg.repeats} repeats, as a
  reliability run, not a new pilot result (D54).
\item \texttt{\FRESHRUN} graded every fresh attempt and each family's reference under the \fv{run.conditions.w}
  conditions, the reference checks and the independent labeler, \fv{run.repeats.w} times per pair under harness
  \fv{run.harness}, after a passed pre-flight (D61).
\item \texttt{\NATURALRUN} graded the first sample's accepted families, every pair \nrv{run.repeats} times, after D19's
  labels and a passed pre-flight (protocol~\S20.2; D76, D78; Section~\ref{sec:natural-cases}).
\end{itemize}
In the new runs a pair whose repeats disagree is \emph{unresolved}, counts as a rejection and is reported.

\paragraph{Measures beyond FA and FR.} The protocol also asks for macro rates with their per-family counts; Youden's
$J$, the share of valid candidates passed less the share of invalid candidates passed; the number of families with an
accepted negative; and negatives by type, that is whether a violated requirement was introduced after the first
request (protocol~\S6). Infrastructure outcomes are reported apart. Its secondary measures are the requirement coverage
of the recovered contracts by a first and a second model mapper, with their agreement (the second mapper, and the
first for \texttt{tinydb}, received an extra instruction; D41), agreement between runs, flakiness over repeated gradings,
valid-task yield with rejection reasons, and generation cost, which was not recorded for the verifiers
(Appendix~\ref{app:deviations}).

\paragraph{Reporting rules.} The pilot keeps the reporting of protocol~\S3: cluster-bootstrap percentile intervals over
families (\dv{design.bootstrap.resamples} resamples), which with \plv{run.families.w} families describe variability
rather than evidence of an effect, with exact binomial intervals as a reference scale. New runs follow protocol~\S15.7
and \S16.6:
\begin{itemize}
\item per-family tables are the primary display, with null results stated alongside;
\item each paired comparison has an exact sign-flip p-value over families with its attainable floor;
\item each distinct graded patch counts once;
\item the oracle-gated view, which drops a verifier in each family where it fails its own reference and reports the lost
  yield, is secondary;
\item bootstrap intervals go to an appendix, marked as not inferential below \dv{design.bootstrap.families} families
  (Table~\ref{tab:fresh-boot}).
\end{itemize}
The estimands are fixed in protocol~\S15.4 for fresh attempts (among them the share of \texttt{u1} attempts that violate
a correction requirement, and how C1 and C3 grade them) and in \S15.6 for the natural arm (among them each condition's
acceptance of the pre-review states under the source-evidence label). Each natural run is reported on its own; with
fewer than \dv{design.size.threshold} accepted tasks, its tasks are exploratory case studies with per-family tables
only and no pooled rate, interval or p-value (protocol~\S19.6, \S20.3). No confirmatory study has been run. A draft
precision plan, not adopted, puts the need at about \plv{plan.families} eligible families, with a natural negative
each, to detect a difference in false acceptance of \plv{plan.delta}; against it, the arm's cap of
\dv{design.family.cap} tasks gives a power of \plv{plan.power.k30} (D49).
